\documentclass{amsart}
% For dealing with references we use the comment environment
\usepackage{verbatim}
\newenvironment{reference}{\comment}{\endcomment}
%\newenvironment{reference}{}{}
\newenvironment{slogan}{\comment}{\endcomment}
\newenvironment{history}{\comment}{\endcomment}

% For commutative diagrams we use Xy-pic
\usepackage[all]{xy}

% We use 2cell for 2-commutative diagrams.
\xyoption{2cell}
\UseAllTwocells

% We use multicol for the list of chapters between chapters
\usepackage{multicol}

% This is generally recommended for better output
\usepackage{lmodern}
\usepackage[T1]{fontenc}

% For cross-file-references

% Package for hypertext links:
\usepackage{hyperref}

% For any local file, say "hello.tex" you want to link to please

% Theorem environments.
%
\theoremstyle{plain}
\newtheorem{theorem}[subsection]{Theorem}
\newtheorem{proposition}[subsection]{Proposition}
\newtheorem{lemma}[subsection]{Lemma}

\theoremstyle{definition}
\newtheorem{definition}[subsection]{Definition}
\newtheorem{example}[subsection]{Example}
\newtheorem{exercise}[subsection]{Exercise}
\newtheorem{situation}[subsection]{Situation}

\theoremstyle{remark}
\newtheorem{remark}[subsection]{Remark}
\newtheorem{remarks}[subsection]{Remarks}

\numberwithin{equation}{subsection}

% Macros
%
\def\lim{\mathop{\mathrm{lim}}\nolimits}
\def\colim{\mathop{\mathrm{colim}}\nolimits}
\def\Spec{\mathop{\mathrm{Spec}}}
\def\Hom{\mathop{\mathrm{Hom}}\nolimits}
\def\Ext{\mathop{\mathrm{Ext}}\nolimits}
\def\SheafHom{\mathop{\mathcal{H}\!\mathit{om}}\nolimits}
\def\SheafExt{\mathop{\mathcal{E}\!\mathit{xt}}\nolimits}
\def\Sch{\mathit{Sch}}
\def\Mor{\mathop{\mathrm{Mor}}\nolimits}
\def\Ob{\mathop{\mathrm{Ob}}\nolimits}
\def\Sh{\mathop{\mathit{Sh}}\nolimits}
\def\NL{\mathop{N\!L}\nolimits}
\def\CH{\mathop{\mathrm{CH}}\nolimits}
\def\proetale{{pro\text{-}\acute{e}tale}}
\def\etale{{\acute{e}tale}}
\def\QCoh{\mathit{QCoh}}
\def\Ker{\mathop{\mathrm{Ker}}}
\def\Im{\mathop{\mathrm{Im}}}
\def\Coker{\mathop{\mathrm{Coker}}}
\def\Coim{\mathop{\mathrm{Coim}}}

% Boxtimes
%
\DeclareMathSymbol{\boxtimes}{\mathbin}{AMSa}{"02}

%
% Macros for moduli stacks/spaces
%
\def\QCohstack{\mathcal{QC}\!\mathit{oh}}
\def\Cohstack{\mathcal{C}\!\mathit{oh}}
\def\Spacesstack{\mathcal{S}\!\mathit{paces}}
\def\Quotfunctor{\mathrm{Quot}}
\def\Hilbfunctor{\mathrm{Hilb}}
\def\Curvesstack{\mathcal{C}\!\mathit{urves}}
\def\Polarizedstack{\mathcal{P}\!\mathit{olarized}}
\def\Complexesstack{\mathcal{C}\!\mathit{omplexes}}
% \Pic is the operator that assigns to X its picard group, usage \Pic(X)
% \Picardstack_{X/B} denotes the Picard stack of X over B
% \Picardfunctor_{X/B} denotes the Picard functor of X over B
\def\Pic{\mathop{\mathrm{Pic}}\nolimits}
\def\Picardstack{\mathcal{P}\!\mathit{ic}}
\def\Picardfunctor{\mathrm{Pic}}
\def\Deformationcategory{\mathcal{D}\!\mathit{ef}}

\setlength{\emergencystretch}{3em}
\begin{document}
\title[Stacks Project, Tag 0H40]{485 --- Cokernels of finite modules over two finite-type algebras become flat over a localized Noetherian domain}
\maketitle
\noindent Copyright (C) 2005--2025 Johan de Jong. Stacks Project authors. Source: \href{https://stacks.math.columbia.edu/tag/0H40}{Tag 0H40}.

\noindent Editorial difficulty note (not author text). Difficulty H2: The cokernel need not be finite over either algebra in the usable form, so ordinary generic flatness alone is insufficient. A prime filtration reduces to a domain, generic freeness embeds the finite source in a finite free summand after localization, and control of the entire localization quotient establishes universal injectivity..

\noindent Editorial provenance note (not author text). History: excerpt from revision \texttt{a0444\allowbreak 6e57e\allowbreak c1fbc\allowbreak 252a8\allowbreak 71afc\allowbreak ec775\allowbreak 2fb28\allowbreak 07b14}, more-morphisms.tex, lines 15700--15761. Reference presentation was adapted; original-author context and core support blocks were retained or added. The mathematical wording is unchanged. Licensed under GNU FDL 1.2 or later; no invariant sections or cover texts. See ../licenses/COPYING. Context blocks reproduce author definitions and supporting results; they are not additional counted problems.

\begin{lemma}
\label{lemma-cokernel-flat}
Let $R$ be a Noetherian domain. Let $R \to A \to B$ be finite type ring maps.
Let $M$ be a finite $A$-module and let $N$ a finite $B$-module.
Let $M \to N$ be an $A$-linear map. There exists an nonzero $f \in R$
such that the cokernel of $M_f \to N_f$ is a flat $R_f$-module.
\end{lemma}

\begin{proof}
By replacing $M$ by the image of $M \to N$, we may assume $M \subset N$.
Choose a filtration $0 = N_0 \subset N_1 \subset \ldots \subset N_t = N$
such that $N_i/N_{i - 1} = B/\mathfrak q_i$ for some prime ideal
$\mathfrak q_i \subset B$, see
Algebra, Lemma \href{https://stacks.math.columbia.edu/tag/00L0}{lemma filter Noetherian module (Tag 00L0)}.
Set $M_i = M \cap N_i$. Then $Q = N/M$ has a filtration by the submodules
$Q_i = N_i/M_i$. It suffices to prove $Q_i/Q_{i - 1}$ becomes flat
after localizing at a nonzero element of $f$ (since extensions of
flat modules are flat by Algebra, Lemma \href{https://stacks.math.columbia.edu/tag/00HM}{lemma flat ses (Tag 00HM)}).
Since $Q_i/Q_{i - 1}$ is isomorphic to the  cokernel
of the map $M_i/M_{i - 1} \to N_i/N_{i - 1}$, we reduce to the
case discussed in the next paragraph.

\medskip\noindent
Assume $B$ is a domain and $M \subset N = B$. After replacing $A$ by the
image of $A$ in $B$ we may assume $A \subset B$. By generic flatness,
we may assume $A$ and $B$ are flat over $R$
(Algebra, Lemma \href{https://stacks.math.columbia.edu/tag/051R}{lemma generic flatness Noetherian (Tag 051R)}).
It now suffices to show $M \to B$ becomes $R$-universally
injective after replacing $R$ by a principal localization
(Algebra, Lemma \href{https://stacks.math.columbia.edu/tag/058P}{lemma ui flat domain (Tag 058P)}).
By generic freeness, we can find a nonzero $g \in A$ such that
$B_g$ is a free $A_g$-module
(Algebra, Lemma \href{https://stacks.math.columbia.edu/tag/051R}{lemma generic flatness Noetherian (Tag 051R)}).
Thus we may choose a direct summand $M' \subset B_g$ as an $A_g$-module,
which is finite free as an $A_g$-module, and
such that $M \to B \to B_g$ factors through $M'$.
Clearly, it suffices to show that $M \to M'$
becomes $R$-universally injective after replacing
$R$ by a principal localization.

\medskip\noindent
Say $M' = A_g^{\oplus n}$. Since $M \subset M'$ is a finite $A$-module,
we see that $M$ is contained in $(1/g^m)A^{\oplus n}$ for some $m \geq 0$.
After changing our basis for $M'$ we may assume $M \subset A^{\oplus n}$.
Then it suffices to show that $A^{\oplus n}/M$ and $A_g/A$ become
$R$-flat after replacing $R$ by a principal localization. Namely, then
$M \to A^{\oplus n}$ and $A^{\oplus n} \to A_g^{\oplus n}$ are
universally injective by Algebra, Lemma \href{https://stacks.math.columbia.edu/tag/00HL}{lemma flat tor zero (Tag 00HL)}
and consequently so is the composition $M \to M' = A_g^{\oplus n}$.

\medskip\noindent
By generic flatness (see reference above), we may assume the
module $A^{\oplus n}/M$ is $R$-flat. For the quotient $A_g/A$
we use the fact that
$$
A_g/A = \colim (1/g^m)A/A \cong \colim A/g^mA
$$
and the module $A/g^mA$ has a filtration of length $m$ whose
successive quotients are isomorphic to $A/gA$. Again by generic
flatness we may assume $A/gA$ is $R$-flat and hence each $A/g^mA$
is $R$-flat, and hence so is $A_g/A$.
\end{proof}
\end{document}
